% UNIFIED 2026-05-13
% SOURCE MAP for Manuscript/plan2_levelset_identity.tex
% ==========================================================================
% Agent 9 deliverable: first subsection of 03_plan2_draft.tex.
%
% Each numbered item below corresponds to a labelled result inside this file
% (see \label{} keys). Status keys:
%   PROVED      : proved in this subsection (full proof written here).
%   ADAPTED     : adapted from a named local building block.
%   CITED       : cited from the literature, with exact theorem number.
%   IMPORT      : imported from preliminaries (Agents 2-4) without reproof.
%
% Items
% -----
% [E1.0]  Notation, standing hypotheses, Talenti L^infty bound.
%           Source: WIP_LevelSetIdentity.tex (Notation+Standing);
%                   00_notation_register.md sections 1-9.
%           Status: ADAPTED (notation register); CITED(Talenti1976).
% [E1.1]  Lemma lem:coarea: -m'(t) = alpha(t).
%           Source: WIP_LevelSetIdentity.tex Lem 2.1; agent1-finite-perimeter-identity.md Lem 1.
%           Status: PROVED. Cited inputs: AFP2000 Thm 3.40, Maggi2012 Thm 13.1.
% [E1.2]  Lemma lem:flux: beta(t) = m(t)  (weak divergence-flux identity).
%           Source: WIP_LevelSetIdentity.tex Lem 3.1; agent1-finite-perimeter-identity.md Lem 2.
%           Status: PROVED. Cited inputs: AFP2000 Thm 3.99, Maggi2012 Thm 13.1,
%                   Maggi2012 Ch.17 (rigorous divergence on finite-perimeter sets).
% [E1.3]  Lemma lem:Iab: layer-cake form of I, AC of I, I' = u^*.
%           Source: WIP_LevelSetIdentity.tex Lem 4.1-4.2; agent1 Lem 3a-3b.
%           Status: PROVED. Cited inputs: Kawohl1985 Sec 2; Talenti1976.
% [E1.4]  Lemma lem:Ipp (corrected I''): I''(s) = -1/alpha(t(s)).
%           Source: WIP_LevelSetIdentity.tex Lem 4.3 (corrected form); agent1 Lem 3c.
%           Status: PROVED. Cited inputs: AFP2000 Thm 3.99 (BV chain rule);
%                   BBC1975 (inverse-function theorem for monotone AC).
% [E1.5]  Lemma lem:Gpp: pointwise convexity formula G''(s) = 1/(c_n s^{1-2/n}) - 1/alpha(t).
%           Status: PROVED (direct).
% [E1.6]  Definition def:DHDI: D_H(t), D_I(t); nonnegativity (Cauchy-Schwarz, De Giorgi).
%           Status: PROVED. Cited inputs: DeGiorgi1958, Maggi2012 Thm 14.1.
% [E1.7]  Lemma lem:cov: pointwise-to-integrated convexity identity.
%           Status: PROVED.
% [E1.8]  Lemma lem:Gamma: Gamma(Omega) = (1/L) int s G''(s) ds.
%           Status: PROVED (integration by parts).
% [E1.9]  Lemma lem:endpoint: Gamma(Omega) = (2/|Omega|)(E(Omega) - E(B^*)).
%           Status: PROVED. Cited inputs: PolyaSzego1951, Bandle1980 II.5.
% [E1.10] Theorem thm:deficit (EXACT DEFICIT IDENTITY).
%           Source: WIP_LevelSetIdentity.tex Thm 5.1 (main theorem).
%           Status: PROVED. Notation register Section 8.
% [E1.11] Corollary cor:variance: weighted variance form of D_H.
%           Source: WIP_LevelSetIdentity.tex Sec 6 Rem; level-set-deficit-identity.md Sec 6.
%           Status: PROVED.
% [E1.12] Lemma lem:profile-gap: profile gap b - v <= 2 delta_T / s.
%           Source: level-set-deficit-identity.md Sec 4 (and Lem 8.1).
%           Status: PROVED.
% [E1.13] Lemma lem:slab-good: existence of a good isoperimetric-defect level
%           in a slab (Lemma 8.2 of source note).
%           Status: PROVED.
% [E1.14] Lemma lem:superlevel-transfer: Lemma 8.3 (superlevel asymmetry transfer).
%           Source: level-set-deficit-identity.md Lem 8.3.
%           Status: PROVED.
%
% Hypothesis tracking
% -------------------
% Standing class of Theorem thm:deficit: Omega subset R^n open, |Omega| < infty.
% No boundedness, smoothness, or connectedness of Omega is required.
% All boundary integrals are on De Giorgi reduced boundaries.
%
% Constants
% ---------
% All constants are explicit and depend only on the dimension n. The constant
% c_n := n^2 omega_n^{2/n} appears throughout. No constants of the form
% C(n,R,rho_*) appear in this subsection: bounded-class constants enter only in
% the downstream Plan 2 blocks (MetricFramework onward), which consume this
% subsection.
%
% Notation register compliance
% ----------------------------
% Dimension n (not N). B^* = ball of volume |Omega|, centred at origin.
% R_Omega := (|Omega|/omega_n)^{1/n} is used in place of the symbol R for the
% volume radius (notation register Section 2, R/R_Omega conflict resolution).
% alpha(t), beta(t), D_H, D_I match notation register Section 8 exactly.
% ==========================================================================

\subsection{Exact level-set deficit identity}
\label{subsec:plan2-levelset-identity}

This subsection contains the analytic foundation of Plan~2. We prove an
exact representation of the (normalised) Saint--Venant deficit
$\delta_{T}(\Omega)$ as a weighted integral, over the levels of the
torsion function, of two nonnegative pointwise defects: a
Cauchy--Schwarz defect $D_{H}(t)$ of $|\nabla u|$ on the reduced
boundary $\partial^{*}E_{t}$, and a squared-perimeter isoperimetric
defect $D_{I}(t)$ of the superlevel set $E_{t}$. The identity is
valid, in full finite-perimeter generality, for any open
$\Omega\subset\R^{n}$ with $0<|\Omega|<\infty$; no smoothness,
boundedness, or connectedness of $\Omega$ is assumed.

The result is the analytic input of every downstream block in
\S\ref{sec:plan2-metric}--\S\ref{sec:plan2-assembly}. Two further
consequences proved here (the profile-gap estimate and the
superlevel-transfer estimate) are used in
\S\ref{sec:plan2-weighted-trace} (good levels in a near-boundary slab)
and in \S\ref{sec:plan2-boundary-layer} (transfer of the Fraenkel
asymmetry from a superlevel set back to $\Omega$).

\subsubsection*{Standing hypotheses and notation}
\label{sssec:lsi-notation}

Throughout the subsection, $n\ge 2$ and $\Omega\subset\R^{n}$ is an open
set with $0<|\Omega|<\infty$. We use the canonical notation register
of \S\ref{sec:notation} without further comment, and recall only the
items that enter the present argument:

\begin{itemize}
\item $u=u_{\Omega}\in H^{1}_{0}(\Omega)$ is the variational torsion
function, i.e.\ the unique minimiser of
$J(v):=\tfrac12\int_{\Omega}|\nabla v|^{2}-\int_{\Omega}v$ on
$H^{1}_{0}(\Omega)$, equivalently the weak solution of
$-\Delta u=1$ in $\Omega$ with zero $H^{1}_{0}$-trace; testing $-\Delta
u=1$ against $u^{-}\in H^{1}_{0}(\Omega)$ gives $u\ge 0$ a.e.

\item $E(\Omega):=-\tfrac12\int_{\Omega}u\,dx=-\tfrac12 T(\Omega)$ is the
BDV-normalised Saint--Venant energy (notation register \S 3), and
$B^{*}\subset\R^{n}$ is the open ball with $|B^{*}|=|\Omega|$, centred at
the origin, of (volume) radius
\[
   R_{\Omega}:=\Bigl(\frac{|\Omega|}{\omega_{n}}\Bigr)^{1/n}.
\]
(We write $R_{\Omega}$ rather than $R$ to avoid the conflict with the
confining radius of \S\ref{sec:notation}.)

\item For $t\ge 0$, $E_{t}:=\{u>t\}$ (precise representative
\cite[\S 15]{Maggi2012}), $m(t):=|E_{t}|$ is right-continuous and
nonincreasing, with $m(0^{+})\le|\Omega|$ and $m(\|u\|_{\infty})=0$. The
decreasing rearrangement is $u^{*}(s):=\inf\{t\ge 0:m(t)\le s\}$ for
$s\in[0,|\Omega|]$.

\item $\partial^{*}E_{t}$ is the De Giorgi reduced boundary
\cite{DeGiorgi1958}, and $\nu_{E_{t}}$ is the measure-theoretic outer
unit normal. We set
\begin{equation}\label{eq:lsi-alpha-beta}
   \alpha(t):=\int_{\partial^{*}E_{t}}\frac{1}{|\nabla u|}\,d\mathcal H^{n-1},
   \qquad
   \beta(t):=\int_{\partial^{*}E_{t}}|\nabla u|\,d\mathcal H^{n-1},
\end{equation}
defined for a.e.\ $t$ (justified by Lemma~\ref{lem:coarea}--\ref{lem:flux}
below), and
\[
   c_{n}:=n^{2}\omega_{n}^{2/n}.
\]
\end{itemize}

\paragraph{Talenti $L^{\infty}$ bound.} By Talenti's pointwise
rearrangement comparison \cite[Theorem 1]{Talenti1976}
(see also \cite[Theorem 1.3.2]{Kesavan2006}),
\begin{equation}\label{eq:lsi-talenti}
   u^{*}_{\Omega}(s)\le u^{*}_{B^{*}}(s)
   =\frac{R_{\Omega}^{2}-(s/\omega_{n})^{2/n}}{2n},
   \qquad s\in[0,|\Omega|],
\end{equation}
so
\begin{equation}\label{eq:lsi-Linfty}
   \|u\|_{L^{\infty}(\Omega)}\le \frac{R_{\Omega}^{2}}{2n}<\infty.
\end{equation}
This is the sole reason boundedness of $\Omega$ is unnecessary at the
analytic level. Combined with $u\in H^{1}_{0}(\Omega)$ extended by zero,
\eqref{eq:lsi-Linfty} gives $u\in BV_{\rm loc}(\R^{n})\cap L^{\infty}$;
the BV coarea formula \cite[Theorem 3.40]{AFP2000},
\cite[Theorem 13.1]{Maggi2012} then applies, and for every Borel
$\varphi\ge 0$,
\begin{equation}\label{eq:lsi-coarea}
   \int_{\Omega}|\nabla u|\,\varphi\,dx
   =\int_{-\infty}^{\infty}\Bigl(\int_{\partial^{*}E_{t}}\varphi\,d\mathcal
   H^{n-1}\Bigr)\,dt.
\end{equation}
The set
\[
   \mathcal R:=\{t\in(0,\|u\|_{\infty}):
   P(E_{t})<\infty\hbox{ and }\partial^{*}E_{t}\subset\Omega\}
\]
has full Lebesgue measure in $(0,\|u\|_{\infty})$, and all statements
``for a.e.\ $t$'' below mean for a.e.\ $t\in\mathcal R$. The countable set
of plateau levels $\{t:|\{u=t\}|>0\}$ is, in particular, Lebesgue-null and
is automatically discarded by the $dt$-integral.

\subsubsection{Step 1: coarea identity for $-m'(t)$}
\label{sssec:lsi-coarea}

\begin{lemma}[Coarea form of $-m'$]\label{lem:coarea}
For a.e.\ $t\in(0,\|u\|_{\infty})$,
\begin{equation}\label{eq:lsi-L1}
   -m'(t)=\alpha(t)
         =\int_{\partial^{*}E_{t}}\frac{1}{|\nabla u|}\,d\mathcal H^{n-1}.
\end{equation}
In particular $m$ is absolutely continuous on $[\eps,\|u\|_{\infty}]$ for
every $\eps>0$, and $\alpha\in L^{1}_{\rm loc}((0,\|u\|_{\infty}))$.
\end{lemma}

\begin{proof}
Fix $0<\alpha_{0}<\beta_{0}<\|u\|_{\infty}$ and apply
\eqref{eq:lsi-coarea} to the Borel function
$\varphi(x)=|\nabla u(x)|^{-1}\mathbf 1_{\{|\nabla u|>0\}}(x)
\mathbf 1_{[\alpha_{0},\beta_{0}]}(u(x))$. The product
$|\nabla u|\,\varphi$ equals $\mathbf 1_{\{\alpha_{0}<u<\beta_{0}\}\cap\{|\nabla
u|>0\}}$ almost everywhere; the contribution of $\{|\nabla u|=0\}$ to
the BV coarea slicing vanishes by \cite[3.40(b)]{AFP2000}, so
\begin{equation}\label{eq:lsi-coarea-app}
   |E_{\alpha_{0}}\setminus E_{\beta_{0}}|
   =\int_{\Omega}\mathbf 1_{\{\alpha_{0}<u<\beta_{0}\}}\,dx
   =\int_{\alpha_{0}}^{\beta_{0}}\!\!\Bigl(\int_{\partial^{*}E_{t}}
     \frac{1}{|\nabla u|}\,d\mathcal H^{n-1}\Bigr)\,dt
   =\int_{\alpha_{0}}^{\beta_{0}}\alpha(t)\,dt.
\end{equation}
Outside the at most countable set of plateau levels,
$|E_{\alpha_{0}}\setminus E_{\beta_{0}}|=m(\alpha_{0})-m(\beta_{0})$, so
$m(\alpha_{0})-m(\beta_{0})=\int_{\alpha_{0}}^{\beta_{0}}\alpha(t)\,dt$.
Differentiating in $\beta_{0}$ at Lebesgue points of $\alpha$ yields
\eqref{eq:lsi-L1}. Local integrability of $\alpha$ on $(0,\|u\|_{\infty})$
follows because \eqref{eq:lsi-coarea-app} bounds
$\int_{\alpha_{0}}^{\beta_{0}}\alpha(t)\,dt\le |\Omega|<\infty$, and
absolute continuity of $m$ on $[\eps,\|u\|_{\infty}]$ follows directly
from \eqref{eq:lsi-coarea-app}.
\end{proof}

\begin{remark}\label{rem:plateaus}
The countable set of plateau levels $\{t:|\{u=t\}|>0\}$ contributes
jumps to $m$ that are absorbed by the $dt$-integral in
Theorem~\ref{thm:deficit} below. Equation~\eqref{eq:lsi-L1} is not
asserted on these levels.
\end{remark}

\subsubsection{Step 2: weak divergence/flux identity}
\label{sssec:lsi-flux}

The next lemma is the rigorous BV/Lipschitz-truncation replacement for
the formal divergence identity
$\int_{\partial^{*}E_{t}}\nabla u\cdot\nu_{E_{t}}\,d\mathcal H^{n-1}
=\int_{E_{t}}\Delta u\,dx$ on the finite-perimeter superlevel $E_{t}$.
The point is that the trace of $\nabla u$ on the reduced boundary
$\partial^{*}E_{t}$ does not exist in general; the Lipschitz truncation
of $u-t$ avoids invoking that trace.

\begin{lemma}[Flux equals volume]\label{lem:flux}
For a.e.\ $t\in(0,\|u\|_{\infty})$, $E_{t}$ has finite perimeter and
\begin{equation}\label{eq:lsi-L2}
   \beta(t)
   =\int_{\partial^{*}E_{t}}|\nabla u|\,d\mathcal H^{n-1}
   =m(t).
\end{equation}
\end{lemma}

\begin{proof}
Fix $t\in\mathcal R$ with $|E_{t}|<\infty$, $P(E_{t})<\infty$, and $u$
not having a plateau at $t$; this holds for a.e.\ $t$.

\smallskip\noindent
\emph{Step 1 (Lipschitz truncation).}
For each $\eps>0$ define
$\phi_{\eps}\colon\R\to[0,1]$ by
$\phi_{\eps}(s):=\min(1,\max(0,s/\eps))$ and set
$\eta_{\eps}(x):=\phi_{\eps}(u(x)-t)$. Since $\phi_{\eps}$ is Lipschitz
with $\phi_{\eps}(0)=0$ and $u\in H^{1}_{0}(\Omega)$, the chain rule for
Sobolev compositions \cite[Theorem 3.99]{AFP2000} yields
$\eta_{\eps}\in H^{1}_{0}(\Omega)$ with
\begin{equation}\label{eq:lsi-eta-grad}
   \nabla\eta_{\eps}(x)
   =\eps^{-1}\nabla u(x)\,\mathbf 1_{\{t<u<t+\eps\}}(x).
\end{equation}
Test the weak equation $-\Delta u=1$ on $H^{1}_{0}(\Omega)$ against
$\eta_{\eps}$:
\begin{equation}\label{eq:lsi-weakform}
   \int_{\Omega}\nabla u\cdot\nabla\eta_{\eps}\,dx
   =\int_{\Omega}\eta_{\eps}\,dx.
\end{equation}
By \eqref{eq:lsi-eta-grad}, the left-hand side of \eqref{eq:lsi-weakform}
equals
\begin{equation}\label{eq:lsi-LHS}
   \int_{\Omega}\nabla u\cdot\nabla\eta_{\eps}\,dx
   =\frac{1}{\eps}\int_{\{t<u<t+\eps\}}|\nabla u|^{2}\,dx.
\end{equation}

\smallskip\noindent
\emph{Step 2 (Coarea applied to the LHS).}
Applying \eqref{eq:lsi-coarea} with
$\varphi(x)=|\nabla u(x)|\,\mathbf 1_{(t,t+\eps)}(u(x))$,
\begin{equation}\label{eq:lsi-RHS-coarea}
   \int_{\{t<u<t+\eps\}}|\nabla u|^{2}\,dx
   =\int_{t}^{t+\eps}\Bigl(\int_{\partial^{*}E_{\tau}}
     |\nabla u|\,d\mathcal H^{n-1}\Bigr)\,d\tau
   =\int_{t}^{t+\eps}\beta(\tau)\,d\tau.
\end{equation}

\smallskip\noindent
\emph{Step 3 (Limit of the right-hand side of \eqref{eq:lsi-weakform}).}
Since $t$ is not a plateau level, $|\{u=t\}|=0$, so
$\eta_{\eps}\to\mathbf 1_{E_{t}}$ pointwise a.e.\ as $\eps\downarrow 0$.
Because $0\le\eta_{\eps}\le \mathbf 1_{\{u>t\}}\le \mathbf 1_{E_{t}}+\mathbf
1_{\{u=t\}}$ and $|E_{t}|<\infty$, the dominated convergence theorem
gives
\[
   \int_{\Omega}\eta_{\eps}\,dx
   \xrightarrow[\eps\downarrow 0]{}|E_{t}|=m(t).
\]

\smallskip\noindent
\emph{Step 4 (Limit of \eqref{eq:lsi-RHS-coarea}).}
The function $\tau\mapsto\beta(\tau)$ is locally integrable on
$(0,\|u\|_{\infty})$: by \eqref{eq:lsi-coarea} with
$\varphi(x)=|\nabla u(x)|\mathbf 1_{(0,\|u\|_{\infty})}(u(x))$, and using
the Dirichlet identity
$\int_{\Omega}|\nabla u|^{2}\,dx=\int_{\Omega}u\,dx<\infty$ (test
$-\Delta u=1$ against $u\in H^{1}_{0}(\Omega)$),
\[
   \int_{0}^{\|u\|_{\infty}}\beta(\tau)\,d\tau
   =\int_{\Omega}|\nabla u|^{2}\,dx
   =\int_{\Omega}u\,dx<\infty.
\]
Therefore, at every Lebesgue point of $\beta$ (i.e.\ for a.e.\ $t$),
\[
   \frac{1}{\eps}\int_{t}^{t+\eps}\beta(\tau)\,d\tau
   \xrightarrow[\eps\downarrow 0]{}\beta(t).
\]

\smallskip\noindent
\emph{Step 5 (Conclusion).}
Combining \eqref{eq:lsi-weakform}--\eqref{eq:lsi-RHS-coarea} with Steps 3
and 4 gives $\beta(t)=m(t)$ for a.e.\ $t\in\mathcal R$.
\end{proof}

\begin{remark}[Formal divergence-theorem reading]\label{rem:divergence}
If one had a one-sided trace of $\nabla u$ on $\partial^{*}E_{t}$ (i.e.\
classical smoothness near the level surface), the relation
$\nabla u\cdot\nu_{E_{t}}=-|\nabla u|$ on $\partial^{*}E_{t}$ would
follow from $u>t$ inside $E_{t}$, and Lemma~\ref{lem:flux} would reduce
to the standard divergence theorem
\cite[Theorem 16.3]{Maggi2012} applied to $\nabla u$:
\[
   m(t)=\int_{E_{t}}1\,dx=-\int_{E_{t}}\Delta u\,dx
   =-\int_{\partial^{*}E_{t}}\nabla u\cdot\nu_{E_{t}}\,d\mathcal H^{n-1}
   =\beta(t).
\]
The Lipschitz-truncation step above is the rigorous replacement when no
such trace is available; see \cite[Chapter 17]{Maggi2012} for the
general theory of fluxes on finite-perimeter sets without traces.
\end{remark}

\subsubsection{Step 3: distributional derivatives of the rearranged
primitive}
\label{sssec:lsi-Iprime}

Define
\begin{equation}\label{eq:lsi-Idef}
   I(s):=\int_{\{u>u^{*}(s)\}}u\,dx,\qquad s\in[0,|\Omega|].
\end{equation}

\begin{lemma}[Layer-cake form of $I$ and first derivative]\label{lem:Iab}
For every $s\in[0,|\Omega|]$,
\begin{equation}\label{eq:lsi-Ilayer}
   I(s)=\int_{0}^{s}u^{*}(\sigma)\,d\sigma,
\end{equation}
$I$ is absolutely continuous on $[0,|\Omega|]$, and $I'(s)=u^{*}(s)$ for
a.e.\ $s\in(0,|\Omega|)$.
\end{lemma}

\begin{proof}
By the equimeasurability of $u$ and $u^{*}$ on $(0,|\Omega|)$
\cite[\S 2]{KawohlBook}, \cite{Talenti1976}, one has
$\int_{\Omega}g(u)\,dx=\int_{0}^{|\Omega|}g(u^{*}(\sigma))\,d\sigma$ for
every Borel $g\ge 0$. Choose $g(t):=t\,\mathbf 1_{\{t>u^{*}(s)\}}$. By
the definition of $u^{*}$, the set $\{\sigma\in(0,|\Omega|):
u^{*}(\sigma)>u^{*}(s)\}$ equals $(0,s)$ up to Lebesgue-null sets, so
\[
   I(s)=\int_{\Omega}g(u)\,dx
       =\int_{0}^{|\Omega|}g(u^{*}(\sigma))\,d\sigma
       =\int_{0}^{s}u^{*}(\sigma)\,d\sigma,
\]
which is \eqref{eq:lsi-Ilayer}. By \eqref{eq:lsi-talenti}, $u^{*}\in
L^{\infty}(0,|\Omega|)\subset L^{1}(0,|\Omega|)$, so \eqref{eq:lsi-Ilayer}
expresses $I$ as the indefinite integral of a bounded Borel function;
the Lebesgue differentiation theorem then gives the absolute continuity
of $I$ on $[0,|\Omega|]$ and $I'(s)=u^{*}(s)$ a.e.
\end{proof}

The next lemma is the {\it corrected} form of the second derivative of
$I$. A naive derivation gives
$I''(s)=-1/a(u^{*}(s))$ with $a(t)=\int_{\partial^{*}E_{t}}|\nabla u|$,
i.e.\ with the {\it flux} $\beta(t)$ in place of the {\it coarea}
$\alpha(t)$. By Lemma~\ref{lem:flux}, $\beta(t)=m(t)$, so that
expression would give $I''(s)=-1/m(u^{*}(s))=-1/s$, which is correct
only when $|\nabla u|$ is $\mathcal H^{n-1}$-a.e.\ constant on
$\partial^{*}E_{t}$ (i.e.\ on the ball). The general identity uses
$\alpha$, and the difference $\alpha\beta-P(E_{t})^{2}$ between the two
is precisely the Cauchy--Schwarz defect $D_{H}$ of \S\ref{sssec:lsi-DHDI}
below.

\begin{lemma}[Corrected second derivative of $I$]\label{lem:Ipp}
The rearrangement $u^{*}$ is absolutely continuous on $[s_{0},|\Omega|]$
for every $0<s_{0}<|\Omega|$, and
\begin{equation}\label{eq:lsi-Ipp}
   I''(s)=(u^{*})'(s)
        =\frac{1}{m'(u^{*}(s))}
        =-\frac{1}{\alpha(u^{*}(s))}
   \qquad\text{for a.e.\ }s\in(0,|\Omega|),
\end{equation}
where $I''$ denotes the absolutely continuous representative of the
distributional derivative of $I'=u^{*}$ on $[s_{0},|\Omega|]$.
\end{lemma}

\begin{proof}
By Lemma~\ref{lem:coarea}, $m$ is absolutely continuous on
$[\eps,\|u\|_{\infty}]$ for every $\eps>0$, with $m'(t)=-\alpha(t)\le 0$
a.e. The set of plateau values $\{t:|\{u=t\}|>0\}$ is at most countable
(monotonicity of $m$); outside it, $m$ is strictly decreasing on its
domain of strict monotonicity. Consequently the right-continuous
generalised inverse $u^{*}$ of $m$ is also strictly decreasing and
absolutely continuous on $[s_{0},|\Omega|]$ for every $s_{0}\in
(0,|\Omega|)$. Furthermore, $\alpha(t)>0$ for a.e.\ $t$ (otherwise $m$
would not be strictly decreasing), so $1/\alpha(t)$ is well-defined
a.e. Applying the inverse-function theorem for monotone absolutely
continuous functions \cite[Theorem 3.99]{AFP2000},
\cite[Lemma 1.4]{BBC1975} to $m(u^{*}(s))=s$ yields
\[
   (u^{*})'(s)=\frac{1}{m'(u^{*}(s))}=-\frac{1}{\alpha(u^{*}(s))}
\]
for a.e.\ $s$. Since $I'(s)=u^{*}(s)$ a.e.\ by Lemma~\ref{lem:Iab}, the
distributional derivative of $I'$ on $[s_{0},|\Omega|]$ coincides with
the classical a.e.\ derivative of $u^{*}$, giving \eqref{eq:lsi-Ipp}.
\end{proof}

\paragraph{Cauchy--Schwarz on the reduced boundary.}
Applying the Cauchy--Schwarz inequality on $\partial^{*}E_{t}$ to the
pair of functions $(|\nabla u|^{1/2},|\nabla u|^{-1/2})$ gives
\begin{equation}\label{eq:lsi-CS}
   P(E_{t})^{2}
   =\Bigl(\int_{\partial^{*}E_{t}}1\,d\mathcal H^{n-1}\Bigr)^{2}
   \le\alpha(t)\,\beta(t),
\end{equation}
with equality if and only if $|\nabla u|$ is $\mathcal H^{n-1}$-a.e.\
constant on $\partial^{*}E_{t}$.

\subsubsection{Step 4: pointwise convexity formula}
\label{sssec:lsi-Gpp}

Following the Nicola--Tilli convexity-gap idea adapted to the torsion
problem, introduce the auxiliary function
\begin{equation}\label{eq:lsi-Gdef}
   G(s):=I(s)+\frac{s^{1+2/n}}{(4+2n)\,\omega_{n}^{2/n}},
   \qquad s\in[0,|\Omega|].
\end{equation}
A direct computation gives
\[
   \frac{d^{2}}{ds^{2}}\frac{s^{1+2/n}}{(4+2n)\omega_{n}^{2/n}}
   =\frac{(1+2/n)(2/n)}{(4+2n)\omega_{n}^{2/n}}s^{2/n-1}
   =\frac{1}{n^{2}\omega_{n}^{2/n}\,s^{1-2/n}}
   =\frac{1}{c_{n}\,s^{1-2/n}}.
\]

\begin{lemma}[Pointwise convexity formula]\label{lem:Gpp}
For a.e.\ $s\in(0,|\Omega|)$, with $t=u^{*}(s)$ (equivalently $s=m(t)$ on
the strictly decreasing branch of $m$),
\begin{equation}\label{eq:lsi-Gpp}
   G''(s)=\frac{1}{c_{n}\,s^{1-2/n}}-\frac{1}{\alpha(t)}.
\end{equation}
\end{lemma}

\begin{proof}
Combine Lemma~\ref{lem:Ipp} ($I''(s)=-1/\alpha(u^{*}(s))$) with the
explicit second derivative of the polynomial corrector
$\frac{s^{1+2/n}}{(4+2n)\omega_{n}^{2/n}}$ computed above.
\end{proof}

\subsubsection{Step 5: the two pointwise defects}
\label{sssec:lsi-DHDI}

\begin{definition}[Cauchy--Schwarz and isoperimetric defects]\label{def:DHDI}
For a.e.\ $t\in(0,\|u\|_{\infty})$, set
\begin{align}
   D_{H}(t)
   &:=\alpha(t)\,\beta(t)-P(E_{t})^{2}
    =\Bigl(\!\int_{\partial^{*}E_{t}}\!\!\tfrac{1}{|\nabla u|}\Bigr)
      \Bigl(\!\int_{\partial^{*}E_{t}}\!\!|\nabla u|\Bigr)-P(E_{t})^{2},
   \label{eq:lsi-DH}\\
   D_{I}(t)
   &:=P(E_{t})^{2}-c_{n}\,m(t)^{2-2/n}.
   \label{eq:lsi-DI}
\end{align}
\end{definition}

\begin{lemma}[Pointwise nonnegativity]\label{lem:DH-DI-pos}\label{lem:lsid:CS-defect}\label{prop:lsid:CS-defect}
For a.e.\ $t\in(0,\|u\|_{\infty})$, $D_{H}(t)\ge 0$ with equality if and
only if $|\nabla u|$ is $\mathcal H^{n-1}$-a.e.\ constant on
$\partial^{*}E_{t}$; and $D_{I}(t)\ge 0$ with equality if and only if
$E_{t}$ is (equivalent to) a ball.
\end{lemma}

\begin{proof}
$D_{H}\ge 0$ is the Cauchy--Schwarz inequality \eqref{eq:lsi-CS}. The
equality case is the standard equality case in Cauchy--Schwarz.

$D_{I}\ge 0$ is the squared form of the De~Giorgi isoperimetric
inequality $P(E_{t})\ge n\omega_{n}^{1/n}|E_{t}|^{1-1/n}$ valid for every
set of finite perimeter \cite[Theorem~II]{DeGiorgi1958},
\cite[Theorem 14.1]{Maggi2012}; equality holds iff $E_{t}$ is
$\mathcal H^{n}$-equivalent to a ball.
\end{proof}

\subsubsection{Step 6: pointwise-to-integrated convexity identity}
\label{sssec:lsi-cov}

\begin{lemma}[Pointwise-to-integrated convexity identity]\label{lem:cov}\label{lem:lsid:DHDI-psi}
With $L:=|\Omega|$,
\begin{equation}\label{eq:lsi-cov}
   \int_{0}^{L}s\,G''(s)\,ds
   =\int_{0}^{\|u\|_{\infty}}\frac{D_{H}(t)+D_{I}(t)}{c_{n}\,m(t)^{1-2/n}}\,dt.
\end{equation}
\end{lemma}

\begin{proof}
By Lemma~\ref{lem:Gpp},
\begin{equation}\label{eq:lsi-cov-A}
   \int_{0}^{L}s\,G''(s)\,ds
   =\int_{0}^{L}s\Bigl(\frac{1}{c_{n}\,s^{1-2/n}}
      -\frac{1}{\alpha(u^{*}(s))}\Bigr)\,ds.
\end{equation}
Perform the change of variable $s=m(t)$, $ds=m'(t)\,dt=-\alpha(t)\,dt$
on the strictly decreasing branch of $m$. The countable set of plateau
levels of $u$ is a Lebesgue-null set in $t$, and contributes nothing to
the integral. As $s$ runs from $0$ to $L$, $t$ runs from $\|u\|_{\infty}$
to $0$; reversing orientation,
\begin{equation}\label{eq:lsi-cov-B}
   \int_{0}^{L}s\,G''(s)\,ds
   =\int_{0}^{\|u\|_{\infty}}m(t)
     \Bigl(\frac{1}{c_{n}\,m(t)^{1-2/n}}-\frac{1}{\alpha(t)}\Bigr)
     \alpha(t)\,dt
   =\int_{0}^{\|u\|_{\infty}}\Bigl(
        \frac{m(t)\,\alpha(t)}{c_{n}\,m(t)^{1-2/n}}-m(t)\Bigr)\,dt.
\end{equation}
By Lemma~\ref{lem:flux}, $m(t)=\beta(t)$ a.e., so
$m(t)\alpha(t)=\alpha(t)\beta(t)$. Adding and subtracting
$P(E_{t})^{2}/(c_{n}m(t)^{1-2/n})$ in the integrand of
\eqref{eq:lsi-cov-B},
\[
   \frac{m(t)\,\alpha(t)}{c_{n}\,m(t)^{1-2/n}}-m(t)
   =\frac{\alpha(t)\,\beta(t)-c_{n}\,m(t)^{2-2/n}}{c_{n}\,m(t)^{1-2/n}}
   =\frac{D_{H}(t)+D_{I}(t)}{c_{n}\,m(t)^{1-2/n}}.
\]
Substituting back into \eqref{eq:lsi-cov-B} gives \eqref{eq:lsi-cov}.
\end{proof}

\subsubsection{Step 7: endpoint convexity gap}
\label{sssec:lsi-Gamma}

With $L:=|\Omega|$, define
\begin{equation}\label{eq:lsi-Gammadef}
   \Gamma(\Omega):=G'(L)-\frac{G(L)}{L}.
\end{equation}

\begin{lemma}[Endpoint convexity-gap formula]\label{lem:Gamma}
$\displaystyle\Gamma(\Omega)=\frac{1}{L}\int_{0}^{L}s\,G''(s)\,ds$.
\end{lemma}

\begin{proof}
$I\in W^{1,\infty}(0,L)\subset W^{2,1}_{\rm loc}((0,L])$ by
Lemmas~\ref{lem:Iab}--\ref{lem:Ipp}, and the polynomial corrector in
\eqref{eq:lsi-Gdef} is smooth on $(0,L]$ and Lipschitz on $[0,L]$; so $G$
is absolutely continuous on $[0,L]$ and in $W^{2,1}_{\rm loc}((0,L])$,
with $G(0)=I(0)=0$. By Lemma~\ref{lem:Iab},
\[
   G'(s)=u^{*}(s)+\frac{(1+2/n)s^{2/n}}{(4+2n)\omega_{n}^{2/n}}
        =u^{*}(s)+\frac{s^{2/n}}{2n\,\omega_{n}^{2/n}},
\]
using $\frac{1+2/n}{4+2n}=\frac{1}{2n}$; $G'$ is bounded on $[0,L]$ by
\eqref{eq:lsi-Linfty}. Integration by parts on $(\eps,L)$ gives
\[
   \int_{\eps}^{L}s\,G''(s)\,ds
   =L\,G'(L)-\eps\,G'(\eps)-G(L)+G(\eps).
\]
As $\eps\downarrow 0$, $\eps\,G'(\eps)\to 0$ and $G(\eps)\to G(0)=0$ by
the boundedness of $G'$ and the absolute continuity of $G$. Therefore
$\int_{0}^{L}s\,G''(s)\,ds=L\,G'(L)-G(L)$; dividing by $L$ yields the
claim.
\end{proof}

\subsubsection{Step 8: endpoint identification with the Saint--Venant
deficit}
\label{sssec:lsi-endpoint}

\begin{lemma}[Endpoint value]\label{lem:endpoint}
\begin{equation}\label{eq:lsi-endpoint}
   \Gamma(\Omega)=\frac{2}{|\Omega|}\bigl(E(\Omega)-E(B^{*})\bigr).
\end{equation}
\end{lemma}

\begin{proof}
At $s=L$, $u^{*}(L)=0$ by definition of $u^{*}$ (the lowest level is the
zero level). By Lemma~\ref{lem:Iab},
$I(L)=\int_{0}^{L}u^{*}(\sigma)\,d\sigma=\int_{\Omega}u\,dx$ (using
equimeasurability), so
\[
   \frac{G(L)}{L}=\frac{1}{L}\int_{\Omega}u\,dx
                  +\frac{L^{2/n}}{(4+2n)\,\omega_{n}^{2/n}}.
\]
Also $I'(L)=u^{*}(L)=0$, so
\[
   G'(L)=\frac{L^{2/n}}{2n\,\omega_{n}^{2/n}}.
\]
Combining,
\begin{equation}\label{eq:lsi-Gamma-explicit}
   \Gamma(\Omega)
   =\frac{L^{2/n}}{2n\,\omega_{n}^{2/n}}
    -\frac{L^{2/n}}{(4+2n)\,\omega_{n}^{2/n}}
    -\frac{1}{L}\int_{\Omega}u\,dx
   =\frac{L^{2/n}}{2n(n+2)\,\omega_{n}^{2/n}}
    -\frac{1}{L}\int_{\Omega}u\,dx,
\end{equation}
where we used $4+2n=2(n+2)$ and
$\frac{1}{2n}-\frac{1}{2(n+2)}=\frac{1}{2n(n+2)}$.

For the ball $B^{*}$ of volume $L=\omega_{n}R_{\Omega}^{n}$, the torsion
function is $u_{B^{*}}(x)=(R_{\Omega}^{2}-|x|^{2})/(2n)$ (the unique
solution of $-\Delta u_{B^{*}}=1$ in $B^{*}$, $u_{B^{*}}=0$ on
$\partial B^{*}$). Direct integration in spherical coordinates,
\[
   \int_{B^{*}}u_{B^{*}}\,dx
   =\frac{1}{2n}\int_{B^{*}}(R_{\Omega}^{2}-|x|^{2})\,dx
   =\frac{n\omega_{n}}{2n}\int_{0}^{R_{\Omega}}(R_{\Omega}^{2}-r^{2})r^{n-1}\,dr
   =\frac{\omega_{n}\,R_{\Omega}^{n+2}}{2n(n+2)}.
\]
Since $L=\omega_{n}R_{\Omega}^{n}$, $L^{2/n}=\omega_{n}^{2/n}R_{\Omega}^{2}$,
hence
\[
   \frac{L^{2/n}}{2n(n+2)\omega_{n}^{2/n}}
   =\frac{R_{\Omega}^{2}}{2n(n+2)}
   =\frac{1}{L}\cdot\frac{\omega_{n}R_{\Omega}^{n+2}}{2n(n+2)}
   =\frac{1}{L}\int_{B^{*}}u_{B^{*}}\,dx.
\]
Inserting back in \eqref{eq:lsi-Gamma-explicit} and using
$E(\Omega)=-\tfrac12\int_{\Omega}u$ (and the analogous identity for $B^{*}$),
\[
   \Gamma(\Omega)
   =\frac{1}{L}\Bigl(\int_{B^{*}}u_{B^{*}}\,dx-\int_{\Omega}u\,dx\Bigr)
   =\frac{2}{L}\bigl(E(\Omega)-E(B^{*})\bigr).
\]
This identification of $\Gamma(B^{*})$ via the explicit ball torsion
function is the classical P\'olya--Szeg\H{o} content
\cite{PolyaSzego1951}, \cite[\S II.5]{Bandle1980}.
\end{proof}

\subsubsection{The exact deficit identity}
\label{sssec:lsi-main}

\begin{theorem}[Exact level-set deficit identity]\label{thm:deficit}\label{thm:lsid:identity}\label{thm:plan2-levelset-id}
Let $n\ge 2$ and let $\Omega\subset\R^{n}$ be open with $0<|\Omega|<\infty$.
Let $u=u_{\Omega}\in H^{1}_{0}(\Omega)$ be the variational torsion
function. Then for a.e.\ $t\in(0,\|u\|_{\infty})$, the superlevel set
$E_{t}=\{u>t\}$ has finite perimeter, the coarea identity
\eqref{eq:lsi-L1}, the flux identity \eqref{eq:lsi-L2}, and the corrected
second-derivative identity \eqref{eq:lsi-Ipp} hold, and
\begin{equation}\label{eq:lsi-deficit}
   \boxed{\;
   \frac{2}{|\Omega|}\bigl(E(\Omega)-E(B^{*})\bigr)
   =\frac{1}{|\Omega|}\int_{0}^{\|u\|_{\infty}}
     \frac{D_{H}(t)+D_{I}(t)}{n^{2}\omega_{n}^{2/n}\,m(t)^{1-2/n}}\,dt,\;}
\end{equation}
where $D_{H},D_{I}$ are the nonnegative pointwise defects of
\eqref{eq:lsi-DH}--\eqref{eq:lsi-DI} and $B^{*}$ is the ball with
$|B^{*}|=|\Omega|$.
\end{theorem}

\begin{proof}
Lemmas~\ref{lem:cov}, \ref{lem:Gamma}, \ref{lem:endpoint} together give
\[
   \frac{2}{|\Omega|}\bigl(E(\Omega)-E(B^{*})\bigr)
   =\Gamma(\Omega)
   =\frac{1}{L}\int_{0}^{L}s\,G''(s)\,ds
   =\frac{1}{L}\int_{0}^{\|u\|_{\infty}}\frac{D_{H}(t)+D_{I}(t)}
     {c_{n}\,m(t)^{1-2/n}}\,dt.
\]
The pointwise statements \eqref{eq:lsi-L1}, \eqref{eq:lsi-L2},
\eqref{eq:lsi-Ipp} are proved in Lemmas~\ref{lem:coarea},
\ref{lem:flux}, \ref{lem:Ipp}.
\end{proof}

The deficit identity~\eqref{eq:lsi-deficit} thus expresses
$\delta_{T}(\Omega)$ in the form
\[
   \delta_{T}(\Omega)
   =E(\Omega)-E(B^{*})
   =\frac{1}{2}\int_{0}^{\|u\|_{\infty}}\bigl(D_{H}(t)+D_{I}(t)\bigr)\,w(t)\,dt,
\]
with the {\it explicit weight}
\begin{equation}\label{eq:lsi-weight}
   w(t):=\frac{1}{c_{n}\,m(t)^{1-2/n}}
        =\frac{1}{n^{2}\omega_{n}^{2/n}\,m(t)^{1-2/n}}.
\end{equation}
This is the form quoted in the master brief and in
\S\ref{sec:notation}; the dimensional constant $c_{n}=n^{2}\omega_{n}^{2/n}$
is the only constant appearing in the identity itself.

\subsubsection*{Cauchy--Schwarz/Serrin variance form of $D_{H}$}

The following corollary recasts $D_{H}$ as a level-set Serrin variance
defect, in the form used in
\S\ref{sec:plan2-weighted-trace}--\S\ref{sec:plan2-boundary-layer} when
combining $D_{H}$ with normal-oscillation bounds.

\begin{corollary}[Weighted variance form of $D_{H}$]\label{cor:variance}
For a.e.\ $t\in(0,\|u\|_{\infty})$, with $f:=|\nabla u|$ on
$\partial^{*}E_{t}$ and $\bar f:=m(t)/P(E_{t})$,
\begin{equation}\label{eq:lsi-variance}
   \int_{\partial^{*}E_{t}}\frac{(f-\bar f)^{2}}{f}\,d\mathcal H^{n-1}
   =\frac{m(t)}{P(E_{t})^{2}}\,D_{H}(t).
\end{equation}
\end{corollary}

\begin{proof}
Direct computation, using Lemma~\ref{lem:flux} ($\int f\,d\mathcal H^{n-1}
=\beta(t)=m(t)$) and $\int 1\,d\mathcal H^{n-1}=P(E_{t})$:
\[
   \int_{\partial^{*}E_{t}}\!\frac{(f-\bar f)^{2}}{f}\,d\mathcal H^{n-1}
   =\int_{\partial^{*}E_{t}}\!\!f\,d\mathcal H^{n-1}
    -2\bar f\,P(E_{t})
    +\bar f^{2}\!\int_{\partial^{*}E_{t}}\!\!\tfrac{1}{f}\,d\mathcal H^{n-1}
   =m(t)-\frac{2 m(t)}{1}+\frac{m(t)^{2}}{P(E_{t})^{2}}\,\alpha(t).
\]
The first two terms collapse to $-m(t)$, and using
$\alpha(t)\beta(t)=\alpha(t)m(t)$,
\[
   \int_{\partial^{*}E_{t}}\!\frac{(f-\bar f)^{2}}{f}\,d\mathcal H^{n-1}
   =\frac{m(t)^{2}}{P(E_{t})^{2}}\alpha(t)-m(t)
   =\frac{m(t)\bigl(\alpha(t)\,m(t)-P(E_{t})^{2}\bigr)}{P(E_{t})^{2}}
   =\frac{m(t)}{P(E_{t})^{2}}\,D_{H}(t),
\]
since $\alpha m=\alpha\beta$. This is \eqref{eq:lsi-variance}.
\end{proof}

\subsubsection{Profile-gap consequences (used in
\S\ref{sec:plan2-weighted-trace}--\S\ref{sec:plan2-boundary-layer})}
\label{sssec:lsi-consequences}

The deficit identity~\eqref{eq:lsi-deficit} controls weighted averages
of $D_{H}+D_{I}$ over levels; the two consequences below extract
{\it pointwise} good levels in a controlled near-boundary slab, together
with a transfer estimate from a superlevel set back to $\Omega$. These
are the only consequences of the deficit identity used in the rest of
Plan~2.

\paragraph{Setup.} Write
\[
   L:=|\Omega|,\quad
   R_{\Omega}:=(L/\omega_{n})^{1/n},\quad
   \delta_{T}(\Omega):=E(\Omega)-E(B^{*}),\quad
   v(s):=u_{\Omega}^{*}(s),\quad
   b(s):=u_{B^{*}}^{*}(s).
\]
By Talenti~\eqref{eq:lsi-talenti}, $b(s)-v(s)\ge 0$ on $[0,L]$, and the
explicit ball profile gives
$b(s)=(R_{\Omega}^{2}-(s/\omega_{n})^{2/n})/(2n)$.

\begin{lemma}[Profile-gap estimate]\label{lem:profile-gap}
The function $h(s):=b(s)-v(s)$ is nonnegative and nonincreasing on
$[0,L]$, $h(L)=0$, and
\begin{equation}\label{eq:lsi-profile-int}
   \int_{0}^{L}h(s)\,ds
   =\int_{B^{*}}u_{B^{*}}\,dx-\int_{\Omega}u\,dx
   =2\,\delta_{T}(\Omega).
\end{equation}
Consequently, for every $s\in(0,L]$,
\begin{equation}\label{eq:lsi-profile-pt}
   0\le b(s)-v(s)\le\frac{2\,\delta_{T}(\Omega)}{s}.
\end{equation}
\end{lemma}

\begin{proof}
By Talenti's level-set differential inequality
$(u_{\Omega}^{*})'(s)\ge(u_{B^{*}}^{*})'(s)$ for a.e.\ $s\in(0,L)$
\cite[Theorem 1, (3.6)]{Talenti1976}, $h=b-v$ has $h'(s)\le 0$ a.e.,
i.e.\ $h$ is nonincreasing. Also $h(L)=b(L)-v(L)=0-0=0$. The
identity~\eqref{eq:lsi-profile-int} follows from $\int_{0}^{L}b=
\int_{B^{*}}u_{B^{*}}$ and $\int_{0}^{L}v=\int_{\Omega}u$
(equimeasurability), combined with
$E(\Omega)-E(B^{*})=\tfrac12(\int_{B^{*}}u_{B^{*}}-\int_{\Omega}u)$.
Since $h$ is nonincreasing and nonnegative on $[0,L]$, for any
$s\in(0,L]$,
\[
   s\,h(s)\le\int_{0}^{s}h(\sigma)\,d\sigma\le\int_{0}^{L}h(\sigma)\,d\sigma
   =2\,\delta_{T}(\Omega),
\]
which is~\eqref{eq:lsi-profile-pt}.
\end{proof}

\paragraph{Good-level slab lemma.} Define the weighted level measure
\begin{equation}\label{eq:lsi-nu}
   d\nu(t):=\frac{dt}{c_{n}\,m(t)^{1-2/n}}.
\end{equation}
Theorem~\ref{thm:deficit} reads
$\int_{0}^{\|u\|_{\infty}}(D_{H}(t)+D_{I}(t))\,d\nu(t)
=2\,\delta_{T}(\Omega)$.

\begin{lemma}[Good level in a near-boundary slab]\label{lem:slab-good}
Let $0<\eta<L/4$, and set
\[
   S_{\eta}:=[L-2\eta,L-\eta],
   \qquad
   T_{\eta}:=v(S_{\eta})=[v(L-\eta),v(L-2\eta)].
\]
There exists a dimensional constant $c=c(n)>0$ such that the following
holds. Assume the smallness condition
\begin{equation}\label{eq:lsi-smallness}
   \delta_{T}(\Omega)\le c\,R_{\Omega}^{2}\,\eta.
\end{equation}
Then for every $s\in S_{\eta}$ the level $t_{s}:=v(s)$ satisfies
\begin{equation}\label{eq:lsi-tslab}
   c\,\frac{R_{\Omega}^{2}}{L}\,\eta
   \le t_{s}\le C\,\frac{R_{\Omega}^{2}}{L}\,\eta,
\end{equation}
for some dimensional $C=C(n)$, and
\begin{equation}\label{eq:lsi-Tslen}
   |T_{\eta}|\ge c\,\frac{R_{\Omega}^{2}}{L}\,\eta.
\end{equation}
Moreover there exists $t_{\eta}^{*}\in T_{\eta}$ with
\begin{equation}\label{eq:lsi-good-level}
   D_{H}(t_{\eta}^{*})+D_{I}(t_{\eta}^{*})
   \le\frac{2\,\delta_{T}(\Omega)}{\nu(T_{\eta})}
   \le C'\,\frac{L^{2-2/n}\,\delta_{T}(\Omega)}{\eta\,R_{\Omega}^{2}}
\end{equation}
for a dimensional constant $C'=C'(n)$, and $\eta\le |\Omega\setminus
E_{t_{\eta}^{*}}|\le 2\eta$.
\end{lemma}

\begin{proof}
{\it Profile-height estimates~\eqref{eq:lsi-tslab}.}
By Lemma~\ref{lem:profile-gap}, on $S_{\eta}\subset[L/2,L]$,
\[
   0\le b(s)-v(s)\le \frac{2\,\delta_{T}(\Omega)}{s}\le \frac{4\,
   \delta_{T}(\Omega)}{L}.
\]
For $s=L-\sigma$, $\sigma\in[\eta,2\eta]$, the explicit ball profile
gives
\[
   b(s)
   =\frac{R_{\Omega}^{2}}{2n}\Bigl(1-(1-\sigma/L)^{2/n}\Bigr).
\]
The elementary one-dimensional inequalities
$\frac{2}{n}\frac{\sigma}{L}\bigl(1-\frac{\sigma}{L}\bigr)^{2/n-1}
\le 1-(1-\sigma/L)^{2/n}\le\frac{2}{n}\frac{\sigma}{L}$ for $0<\sigma<L/4$
give, with new dimensional constants $c_{1},C_{1}$,
\[
   c_{1}\frac{R_{\Omega}^{2}}{L}\sigma\le b(s)\le C_{1}\frac{R_{\Omega}^{2}}{L}\sigma.
\]
Combined with the profile gap and \eqref{eq:lsi-smallness} chosen so
that $4\delta_{T}/L\le c_{1}R_{\Omega}^{2}\eta/(2L)$, this gives
\eqref{eq:lsi-tslab} after relabelling $c$, $C$.

{\it Lower bound~\eqref{eq:lsi-Tslen}.}
$|T_{\eta}|=v(L-2\eta)-v(L-\eta)$. Using
$v=b-h$ and the monotonicity of $h$,
\[
   v(L-2\eta)-v(L-\eta)
   =[b(L-2\eta)-b(L-\eta)]-[h(L-2\eta)-h(L-\eta)]
   \ge[b(L-2\eta)-b(L-\eta)]-\frac{4\delta_{T}(\Omega)}{L},
\]
since $0\le h(L-2\eta)-h(L-\eta)\le h(L-2\eta)\le 4\delta_{T}/L$ (using
$h(L-2\eta)\le 2\delta_{T}/(L-2\eta)\le 4\delta_{T}/L$). On the other
hand,
$b(L-2\eta)-b(L-\eta)\ge c_{2}R_{\Omega}^{2}\eta/L$ for a dimensional
$c_{2}$, by the same elementary one-dimensional estimate above. Under
\eqref{eq:lsi-smallness} (with $c\le c_{2}/8$),
$4\delta_{T}/L\le c_{2}R_{\Omega}^{2}\eta/(2L)$, giving
\eqref{eq:lsi-Tslen}.

{\it Good-level existence~\eqref{eq:lsi-good-level}.}
By the area formula (equivalently the coarea identity~\eqref{eq:lsi-L1}
integrated),
\[
   \nu(T_{\eta})=\int_{T_{\eta}}\frac{dt}{c_{n}m(t)^{1-2/n}}
                \ge\frac{1}{c_{n}L^{1-2/n}}\,|T_{\eta}|
                \ge\frac{c\,R_{\Omega}^{2}\,\eta}{c_{n}\,L^{2-2/n}}
\]
(using $m(t)\le L$ on $T_{\eta}$ and \eqref{eq:lsi-Tslen}). The
deficit identity~\eqref{eq:lsi-deficit} integrated over $T_{\eta}$ and
the mean-value inequality give the existence of $t_{\eta}^{*}\in T_{\eta}$
with
\[
   D_{H}(t_{\eta}^{*})+D_{I}(t_{\eta}^{*})
   \le\frac{1}{\nu(T_{\eta})}\int_{T_{\eta}}(D_{H}+D_{I})\,d\nu
   \le\frac{2\,\delta_{T}(\Omega)}{\nu(T_{\eta})}.
\]
The second inequality in \eqref{eq:lsi-good-level} now follows from the
lower bound on $\nu(T_{\eta})$. Finally, $t_{\eta}^{*}\in T_{\eta}$ means
$m(t_{\eta}^{*})\in[L-2\eta,L-\eta]$, i.e.\ $\eta\le|\Omega\setminus
E_{t_{\eta}^{*}}|\le 2\eta$.
\end{proof}

\paragraph{Superlevel transfer of Fraenkel asymmetry.}
This is the elementary boundary-layer transfer used in
\S\ref{sec:plan2-boundary-layer}. Let
$\Fr(\Omega):=\inf_{x_{0}\in\R^{n}}|\Omega\Delta(B^{*}+x_{0})|/|B^{*}|$
be the Fraenkel asymmetry (notation register \S 5).

\begin{lemma}[Superlevel asymmetry transfer]\label{lem:superlevel-transfer}
Let $E_{t}=\{u>t\}$ with $E_{t}\subset\Omega$, and let
$s:=|E_{t}|=L-\eta$ with $0<\eta<L$. Let $B^{(s)}$ be a ball with
$|B^{(s)}|=s$ realising the Fraenkel infimum for $E_{t}$, and let
$B^{(L)}$ be the concentric ball of volume $L$. Then
\begin{equation}\label{eq:lsi-transfer}
   \Fr(\Omega)\le\Fr(E_{t})+\frac{4\eta}{L}.
\end{equation}
\end{lemma}

\begin{proof}
By the triangle inequality for symmetric differences,
\[
   |\Omega\Delta B^{(L)}|
   \le|\Omega\setminus E_{t}|+|E_{t}\Delta B^{(s)}|+|B^{(s)}\Delta B^{(L)}|.
\]
Now $|\Omega\setminus E_{t}|=\eta$, $|E_{t}\Delta B^{(s)}|=s\,
\Fr(E_{t})\le L\,\Fr(E_{t})$, and $B^{(s)}\subset B^{(L)}$
so $|B^{(s)}\Delta B^{(L)}|=L-s=\eta$. Hence
\[
   \Fr(\Omega)\le\frac{|\Omega\Delta B^{(L)}|}{L}
   \le \frac{|E_{t}\Delta B^{(s)}|}{L}+\frac{2\eta}{L}
   \le\Fr(E_{t})+\frac{2\eta}{L}.
\]
Allowing for the asymmetry to be computed at a different centre than the
optimal one (which can cost at most $|\Omega\setminus E_{t}|/L=\eta/L$)
gives \eqref{eq:lsi-transfer} after enlarging the constant from $2$ to
$4$ (any centre admissible for $\Fr(E_{t})$ is admissible for
$\Fr(\Omega)$ up to the boundary-layer correction).
\end{proof}

\subsubsection{Failure modes and what is not assumed}
\label{sssec:lsi-failure}

We close this subsection by recording the regularity assumptions that
are {\it not} used and the ones that are. Both lists are required
inputs of Agent~16's Plan~2 audit (\S\ref{sec:audit-plan2}).

\begin{enumerate}
\item {\it $\Omega$ need not be bounded.} The only consequence of
$|\Omega|<\infty$ is the Talenti $L^{\infty}$ bound \eqref{eq:lsi-Linfty}.
Boundedness $\Omega\subset B_{R}$ enters only later in
\S\ref{sec:plan2-metric}--\S\ref{sec:plan2-assembly}, where strong
isoperimetric inequalities and radial-trace estimates are used.
\item {\it $\Omega$ need not be smooth.} Lemmas~\ref{lem:coarea} and
\ref{lem:flux} use only BV coarea \cite[Theorem 13.1]{Maggi2012} and the
Lipschitz truncation of $u-t$; neither requires $\partial\Omega$ to be a
manifold. There is no use of a one-sided trace of $\nabla u$ on
$\partial^{*}E_{t}$ (cf.\ Remark~\ref{rem:divergence}).
\item {\it $\Omega$ need not be connected.} The torsion function
decouples across connected components and the identity is additive.
\item {\it The identity is a.e., not pointwise.} Lemmas~\ref{lem:coarea},
\ref{lem:flux}, \ref{lem:Ipp} hold for $t$ outside an $\mathcal L^{1}$-null
set of bad levels: at most countably many plateau levels of $u$, the
Sobolev--Sard exceptional set, and BV-coarea exceptional levels. Levels
in this null set contribute zero to the $dt$-integral.
\item {\it All integrals are on reduced boundaries.} For a.e.\ $t$,
$\mathcal H^{n-1}(\partial E_{t}\setminus\partial^{*}E_{t})=0$
\cite{DeGiorgi1958}, \cite[Theorem 16.2]{Maggi2012}, so cusps,
self-tangencies, and the unreduced part of $\partial E_{t}$ contribute
nothing. No graph parametrisation of $\partial^{*}E_{t}$ is used or
available in general.
\item {\it The corrected $I''$.} The use of $\alpha$ (and not $\beta$) in
Lemma~\ref{lem:Ipp} is essential. The naive guess $I''=-1/\beta=-1/m$
holds only on the ball, where $|\nabla u|$ is level-constant and
Cauchy--Schwarz is sharp; in general it would silently delete $D_{H}$
from \eqref{eq:lsi-deficit}.
\end{enumerate}

\paragraph{Where the identity is used downstream.}
The metric framework (\S\ref{sec:plan2-metric}) extracts a
first-variation/velocity defect from the integrated combination
$D_{H}+D_{I}$ via the homothetic radial vector field. The weighted
metric trace (\S\ref{sec:plan2-weighted-trace}) uses
Lemma~\ref{lem:slab-good} to extract a good isoperimetric-defect level
in a near-boundary slab, then upgrades the integrated action bound to a
pointwise endpoint at $\widehat\rho$. The boundary-layer transfer
(\S\ref{sec:plan2-boundary-layer}) uses
Lemma~\ref{lem:superlevel-transfer} to transfer the asymmetry from
$E_{\widehat\rho}$ back to $\Omega$, squaring the resulting
$O(\sqrt{\delta_{T}})$ boundary layer to preserve the sharp exponent.
